% Partie : méthodes
% Chapitre : induction mathématique
% Section : relations et fonctions

\documentclass[../../../include/open-logic-section]{subfiles}

\begin{document}
\olfileid[fr]{mth}{ind}{rel}

\olsection{Relations et fonctions}

Après avoir défini
inductivement un
ensemble d'objets,
comme les nombres
naturels ou les
termes bien formés,
nous pouvons aussi
définir par induction
des \emph{relations
sur} ces objets.
Par exemple, un
terme bien formé~$t_1$
est un sous-terme
d'un autre terme
bien formé~$t_2$
s'il apparaît
comme l'une
de ses parties.
Notons cette relation
$t_1 \sqsubseteq t_2$.
Tout terme est
bien sûr un
sous-terme de
lui-même :
$t \sqsubseteq t$.
La relation reçoit
la définition
inductive suivante.

\begin{defn}
  La relation
  «~le terme bien
  formé~$t_1$
  est un sous-terme
  de~$t_2$~»,
  notée
  $t_1 \sqsubseteq t_2$,
  est définie
  par induction
  sur~$t_2$ :
  \begin{enumerate}
  \item Si $t_2$
    est une lettre,
    alors
    $t_1 \sqsubseteq t_2$
    si et seulement
    si $t_1=t_2$.
  \item Si
    $t_2=[s_1 \circ s_2]$,
    alors
    $t_1 \sqsubseteq t_2$
    si et seulement
    si $t_1=t_2$,
    $t_1 \sqsubseteq s_1$
    ou
    $t_1 \sqsubseteq s_2$.
  \end{enumerate}
\end{defn}

Cette définition donne,
par exemple,
$\mathrm{a}
\sqsubseteq
[\mathrm{b}\circ\mathrm{a}]$.
En effet, d'après
la clause (2),
$\mathrm{a}\sqsubseteq
[\mathrm{b}\circ\mathrm{a}]$
vaut si et
seulement si
$\mathrm{a}
=[\mathrm{b}\circ\mathrm{a}]$,
ou
$\mathrm{a}
\sqsubseteq\mathrm{b}$,
ou
$\mathrm{a}
\sqsubseteq\mathrm{a}$.
Les deux premières
possibilités sont
fausses : la lettre
$\mathrm{a}$ n'est
pas identique au
terme composé
$[\mathrm{b}\circ\mathrm{a}]$,
et
la clause (1) dit
que
$\mathrm{a}
\sqsubseteq\mathrm{b}$
si et seulement si
$\mathrm{a}=\mathrm{b}$,
ce qui est également
faux. En revanche,
toujours d'après
la clause (1),
$\mathrm{a}
\sqsubseteq\mathrm{a}$
si et seulement si
$\mathrm{a}=\mathrm{a}$,
ce qui est vrai.

Le bon fonctionnement
de cette définition
dépend d'un fait
que nous n'avons
pas encore démontré :
chaque terme bien
formé~$t$ est
soit une lettre
seule, soit de
la forme
$[s_1\circ s_2]$
pour des termes
bien formés
$s_1$ et $s_2$
\emph{déterminés de
manière unique}.
Cette unicité
signifie que si
$t=[s_1\circ s_2]$,
on ne peut pas
avoir \emph{aussi}
$t=[r_1\circ r_2]$
avec $s_1\neq r_1$
ou $s_2\neq r_2$.
Sinon, la clause~(2)
pourrait donner
des réponses
contradictoires :
la décomposition
$[s_1\circ s_2]$
pourrait conduire
à $t_1\sqsubseteq t_2$,
alors que
$[r_1\circ r_2]$
conduirait à la
négation de
$t_1\sqsubseteq t_2$.
Avant
de démontrer que
cela n'arrive pas
pour les termes
bien formés,
considérons un
exemple où cela
\emph{peut} arriver.

\begin{defn}
  Définissons
  inductivement
  les \emph{termes
  sans crochets} :
  \begin{enumerate}
  \item Chaque lettre
    est un terme
    sans crochets.
  \item Si $s_1$
    et $s_2$
    sont des termes
    sans crochets,
    $s_1\circ s_2$
    l'est aussi.
  \item Aucun autre
    objet n'est
    un terme
    sans crochets.
  \end{enumerate}
\end{defn}

Les expressions
$\mathrm{a}$,
$\mathrm{b}\circ\mathrm{d}$
et
$\mathrm{b}\circ\mathrm{a}
\circ\mathrm{b}$
sont, par exemple,
des termes sans
crochets. Si nous
définissions pour
eux «~sous-terme~»
de la même façon
que précédemment,
la seconde clause
serait :
\begin{center}
  Si $t_2=s_1\circ s_2$,
  alors
  $t_1\sqsubseteq t_2$
  si et seulement si
  $t_1=t_2$,
  $t_1\sqsubseteq s_1$
  ou
  $t_1\sqsubseteq s_2$.
\end{center}

Or
$\mathrm{b}\circ\mathrm{a}
\circ\mathrm{b}$
peut se lire
$s_1\circ s_2$
avec
\begin{align*}
  s_1 &= \mathrm{b}
    \text{ et} &
  s_2 &= \mathrm{a}\circ\mathrm{b}.
\intertext{On peut aussi le lire
  $r_1\circ r_2$ avec}
  r_1 &= \mathrm{b}\circ\mathrm{a}
    \text{ et} &
  r_2 &= \mathrm{b}.
\end{align*}
Le terme
$\mathrm{a}\circ\mathrm{b}$
est-il un sous-terme
de
$\mathrm{b}\circ\mathrm{a}
\circ\mathrm{b}$ ?
Oui selon la première
lecture, non selon
la seconde.

La propriété selon
laquelle un terme
bien formé admet
une seule décomposition
en termes bien
formés plus simples
s'appelle
l'\emph{unicité de
la lecture}.
Comme les définitions
inductives de relations
sur de tels objets
sont importantes,
nous devons établir
cette propriété.

\begin{prop}
  Supposons que $t$
  soit un terme bien
  formé. Alors,
  soit $t$ est une
  lettre seule,
  soit il existe
  des termes bien
  formés $s_1$
  et $s_2$
  déterminés de
  manière unique
  tels que
  $t=[s_1\circ s_2]$.
\end{prop}

\begin{proof}
  Si $t$ est une
  lettre seule,
  la condition est
  satisfaite.
  Sinon, la définition
  inductive impose
  la forme
  $t=[s_1\circ s_2]$
  pour certains
  termes bien formés
  $s_1$ et~$s_2$.
  Montrons que
  ces deux termes
  sont déterminés
  de manière unique :
  si
  $t=[r_1\circ r_2]$,
  alors $s_1=r_1$
  et $s_2=r_2$.

  Supposons donc
  $t=[s_1\circ s_2]
  =[r_1\circ r_2]$,
  les quatre
  sous-termes étant
  bien formés.
  Les termes $s_1$
  et $r_1$ doivent
  être identiques :
  sinon, puisque
  les deux lectures
  portent sur la
  même suite de
  symboles, l'un
  serait un préfixe
  propre non vide
  de l'autre.
  Or la
  \olref[sti]{prop:initial}
  l'interdit lorsque
  les deux termes sont
  bien formés.
  %
  %
  %
  %
  %
  %
  %
  %
  %
  Ainsi $s_1=r_1$ ;
  il s'ensuit
  immédiatement
  $s_2=r_2$.
\end{proof}

Nous pouvons aussi
définir des fonctions
inductivement.
Par exemple,
définissons la
fonction~$f$
qui associe à
chaque terme bien
formé la profondeur
maximale de ses
crochets emboîtés
$[\dots]$ :
\begin{defn}
  \ollabel{defn:depth}
  La \emph{profondeur}
  d'un terme
  bien formé~$t$,
  notée $f(t)$,
  est définie
  inductivement par
  \[
    f(t)=
    \begin{cases}
      0
        & \text{si $t$ est une lettre},\\
      \max(f(s_1),f(s_2))+1
        & \text{si $t=[s_1\circ s_2]$.}
    \end{cases}
  \]
\end{defn}

Par exemple :
\begin{align*}
  f([\mathrm{a}\circ\mathrm{b}])
    &=\max(f(\mathrm{a}),f(\mathrm{b}))+1\\
    &=\max(0,0)+1=1,
      \text{ et}\\
  f([[\mathrm{a}\circ\mathrm{b}]
      \circ\mathrm{c}])
    &=\max(f([\mathrm{a}\circ\mathrm{b}]),
           f(\mathrm{c}))+1\\
    &=\max(1,0)+1=2.
\end{align*}

Nous supposons ici
que $s_1$ et $s_2$
sont des termes bien
formés et utilisons
le fait que tout
terme bien formé
est soit une lettre,
soit de la forme
$[s_1\circ s_2]$.
L'unicité de cette
décomposition est
encore essentielle.
Reprenons, pour le
voir, les termes
sans crochets
définis plus haut.
La «~définition~»
analogue serait :
\[
  g(t)=
  \begin{cases}
    0
      & \text{si $t$ est une lettre},\\
    \max(g(s_1),g(s_2))+1
      & \text{si $t=s_1\circ s_2$.}
  \end{cases}
\]
Considérons maintenant
le terme sans crochets
$\mathrm{a}\circ\mathrm{b}
\circ\mathrm{c}\circ\mathrm{d}$.
On peut notamment
le lire
$s_1\circ s_2$
avec
\begin{align*}
  s_1 &= \mathrm{a}
    \text{ et} &
  s_2 &= \mathrm{b}\circ\mathrm{c}
    \circ\mathrm{d},
\intertext{ou $r_1\circ r_2$ avec}
  r_1 &= \mathrm{a}\circ\mathrm{b}
    \text{ et} &
  r_2 &= \mathrm{c}\circ\mathrm{d}.
\end{align*}
Le calcul de $g$
selon la première
lecture donnerait
\begin{align*}
  g(s_1\circ s_2)
    &=\max(g(\mathrm{a}),
      g(\mathrm{b}\circ\mathrm{c}
      \circ\mathrm{d}))+1\\
    &=\max(0,2)+1=3.
\intertext{Selon l'autre lecture,
  on obtient}
  g(r_1\circ r_2)
    &=\max(g(\mathrm{a}\circ\mathrm{b}),
      g(\mathrm{c}\circ\mathrm{d}))+1\\
    &=\max(1,1)+1=2.
\end{align*}
Une fonction doit
cependant attribuer
une valeur unique
à chaque argument.
Notre «~définition~»
de~$g$ ne définit
donc pas de fonction.

\begin{prob}
  Donnez une définition
  inductive de la
  fonction~$l$,
  où $l(t)$
  désigne le
  nombre de symboles
  du terme bien
  formé~$t$.
\end{prob}

\begin{prob}
  Démontrez par
  induction structurelle
  sur les termes
  bien formés~$t$
  que $f(t)<l(t)$,
  où $l(t)$ est
  le nombre de
  symboles de~$t$
  et $f(t)$ sa
  profondeur,
  définie dans la
  \olref[mth][ind][rel]{defn:depth}.
\end{prob}

\end{document}
